%%
%% fall-lecture14.tex
%% 
%% Made by Alex Nelson <pqnelson@gmail.com>
%% Login   <alex@lisp>
%% 
%% Started on  2026-11-03T07:09:05-0800
%% Last update 2026-11-03T07:09:05-0800
%% 

\lecture[Pivot to Thermodynamics]{}

\begin{recall}
We just derived the Entropy Equation~\eqref{eq:fall:lec13:entropy-equation},
\begin{equation}
\rho T\frac{D s}{D t}=\Phi-\partial_{i}q_{i}+\rho Q.
\end{equation}
\end{recall}

\begin{node}
Entropy is usually defined in Thermodynamics using the equation
\begin{equation}
T\,\D S = \D Q,
\end{equation}
and we know $S=S(p,T)$ is a function of pressure $p$ and temperature $T$.
Then we can form its ``difference equation''
\begin{equation}\label{eq:fall:lec14:eq1}
\D S = \left(\frac{\partial S}{\partial p}\right)_{T}\,\D p +
\left(\frac{\partial S}{\partial T}\right)_{p}\,\D T
\end{equation}
where $(\partial\dots/\partial\dots)_{p}$ means we are taking the
partial derivative while holding $p$ constant,
$(\partial\dots/\partial\dots)_{T}$ means we are taking the partial
derivative while holding $T$ constant. This is common notation in Thermodynamics.
\end{node}

\begin{definition}
We can also define the \define{Specific Heat Capacity}
\begin{equation}
c_{p} := T \left(\frac{\partial S}{\partial T}\right)_{p}
\end{equation}
where the partial differentiation is taken while holding pressure $p$ constant.
\end{definition}

\begin{node}[Deriving the temperature equation]
We can show
\begin{subequations}
  \begin{align}
\left(\frac{\partial S}{\partial T}\right)_{p}
&= - \left(\frac{\partial V}{\partial T}\right)_{p}\quad\mbox{by a Maxwell relation}\\
&= -\alpha V\quad\mbox{where $\alpha$ is coefficient of Thermal capacity}\\
&= -\alpha/\rho\quad\mbox{since $1/\rho$ is the specific volume}.
  \end{align}
\end{subequations}
Then plugging this back into Eq~\eqref{eq:fall:lec14:eq1},
\begin{equation}
\D S = \frac{c_{P}}{T}\D T-\frac{\alpha}{\rho}\D p.
\end{equation}
Then ``using material derivatives (by dividing through by $\D t$)''
(which I am highly distressed by this cavalier attitude)
\begin{equation}
\frac{D S}{D t} = \frac{c_{P}}{T}\frac{D T}{D t}-\frac{\alpha}{\rho}\frac{D p}{Dt}.
\end{equation}
Then we plug this back into the initial Entropy equation to obtain:
\begin{equation}
\boxed{\rho c_{p}\frac{DT}{Dt}=\alpha T\frac{Dp}{Dt}+\Phi-\partial_{i}q_{i}+\rho Q.}
\end{equation}
This is called the \define{Temperature Equation}.
\end{node}

\begin{node}
Observe for liquids and gases at low Mach numbers (i..e.,
compressibility is limited),
\begin{equation}
\left\|\alpha T\frac{Dp}{Dt}\right\|\ll\left\|\rho T\frac{DT}{Dt}\right\|,
\end{equation}
so we can eliminate $Dp/Dt$ to simplify the temperature equation
\emph{in this particular situation} as:
\begin{equation}
\rho c_{p}\frac{DT}{Dt}=\Phi-\partial_{i}q_{i}+\rho Q.
\end{equation}
\end{node}

\begin{assume}[Temporarily]
\begin{enumerate}
\item No heat source: $Q=0$
\item Newtonian constitutive relation: $\Phi=\tau_{ij}\partial_{j}u_{i}$
\item Fourier heat Law: $q_{i}=\kappa\partial_{i}T$
\end{enumerate}
\end{assume}

\begin{node}
Under these temporary assumptions, we can simplify the temperature
equation even further:
\begin{subequations}
  \begin{align}
\partial_{t}T+u_{i}\partial_{i}T &= 2\frac{\mu}{\rho c_{p}}S_{ij}S_{ij}+\frac{\kappa}{\rho c_{p}}\partial_{i}^{2}T\\
&=\underbrace{2\frac{\nu}{c_{p}}S_{ij}S_{ij}}_{\text{dissipation}} + \underbrace{D_{T}\partial_{i}^{2}T}_{\text{diffusivity}},
  \end{align}
\end{subequations}
where $D_{T}$ is the \define{Diffusivity Coefficient of $T$}, the
first term describes dissipation, the second term describes diffusivity.

\textsc{Observe}: under our assumptions, this (plus the equations of
motion [i.e., conservation of momentum] and conservation of mass)
forms a solvable system.
\end{node}

\begin{recall}
We introduced~(\S\ref{node:lec13:schmidt-number}) the Schmidt number,
\begin{equation}
\Schmidt = \frac{\nu}{D_{T}}  = \frac{\nu c_{p}\rho}{\kappa}=\frac{\mu c_{p}}{\kappa}=:\Prandtl
\end{equation}
where $\Prandtl$ is the \define{Prandtl Number}. The Prandtl number is
\textbf{specifically for temperature diffusivity}, and for all other
diffusivities we use instead the Schmidt number.
\end{recall}

\begin{node}[Boussinesq approximation]
Typically, density is a function of temperature $\rho=\rho(T)$. But there are
times when gravity is important where we use the \define{Boussinesq Approximation}
[pronounced like ``Bus-in-esk''], where we treat $\rho$ as constant
\emph{except} when appearing in the gravity term where we use
\begin{equation}
\rho(T)=\rho_{0}[1-\alpha(T-T_{0})],
\end{equation}
where $T_{0}$ is the reference temperature, $\rho_{0}$ is the
approximately constant density. Then the conservation of momentum
equations become
\begin{equation}\label{eq:fall:lec14:boussinesq-approx-momentum-eqns}
\rho_{o}\frac{Du_{i}}{Dt}=-\partial_{i}p+\mu\partial_{j}^{2}u_{i}-\rho_{0}[1-\alpha(T-T_{0})]g\delta_{i,3}
\end{equation}
assuming that $-g\widehat{e}_{3}$ is the gravitational force vector.

Why do this? Well, it ``kinda decouples'' the temperature equation
from the momentum equation, in the sense that temperature $T$ enters
the momentum equation~\eqref{eq:fall:lec14:boussinesq-approx-momentum-eqns}
as an \emph{external} body force rather than appearing ``with full
weight'' on the right-hand side of the equation. Numerically, this
allows us to evolve the temperature from $T(t_{n})$ to $T(t_{n+1})$
using the velocity $u_{i}(t_{n})$, then plug this $T(t_{n+1})$ to
determine the velocity at the next time step $u_{i}(t_{n+1})$.
\end{node}

\begin{caution}
The Boussinesq approximation works well if $\rho$ changes very slowly.
\end{caution}

\begin{remark}
People cite Boussinesq's book~\cite{boussinesq1897theorie} for the origin of the Boussinesq
approximation. 
\end{remark}

\subsection{Continuing Thermodynamics}

\begin{puzzle}
How do we convert one set of thermodynamic variables to another?
\end{puzzle}

\begin{law}[First]
The First Law of Thermodynamics gives us
\begin{align}
\D e &= \D Q + \D W\\
\begin{pmatrix}
\mbox{change in}\\
\mbox{stored/internal}\\
\mbox{energy}
\end{pmatrix} &= \begin{pmatrix}\mbox{heat transfer}\\
\mbox{to the system}
\end{pmatrix}
+ \begin{pmatrix}\mbox{work done}\\
\mbox{on the system}
\end{pmatrix}\nonumber
\end{align}
\end{law}

\begin{caution}
Care should be taken because I am being sloppy about so-called
\define{Inexact Differentials} which are typically written using
$\delta Q$ (instead of $\D Q$) and $\delta W$ (instead of $\D W$).
Inexact differentials are \emph{path dependent}. We should really have
\begin{equation}
\D e = \delta Q + \delta W,
\end{equation}
but I do not believe it is too serious a consequence here.
\end{caution}

\begin{node}
In equilibrium processes (we have a slow change), we have the
definition of entropy
\begin{equation}
T\,\D S = \D Q,
\end{equation}
and the work done by the system
\begin{equation}
\D W = -p\,\D V,
\end{equation}
where $V=1/\rho$ is the specific volume. Then we have the
\define{Canonical Equation of State}:
\begin{equation}\label{eq:canonical-eqn-of-state}
\boxed{\D e = T\,\D S - p\,\D V.}
\end{equation}
\end{node}

\begin{node}
We are saying that internal energy is a function of entropy and
specific volume $e=e(S,V)$ but it can also be a function of pressure
and temperature $e=e(p,T)$. We only need two parameters for internal
energy (or any other quantity of interest). If we tried using three
parameters (for example, $e=e(S,V,T)$), then we could parametrize one
of the arguments in terms of the other two $T=T(S,V)$ which reduces
the internal energy to be a function of 2 parameters again.

The upshot: everything in thermodynamics will be a function of two parameters.
\end{node}

\begin{node}[Maxwell relation for energy]
We see that internal energy is a function of entropy and volume $e=e(S,V)$
and also satisfies the canonical Equation of State Eq~\eqref{eq:canonical-eqn-of-state}.
This means we can set the difference equation equal to the canonical
equation of state and match coefficients of the differentials:
\begin{subequations}
  \begin{align}
    \D e &= \left(\frac{\partial e}{\partial S}\right)_{V}\,\D S
    + \left(\frac{\partial e}{\partial V}\right)_{S}\,\D V\\
    &= T\,\D S - p\,\D V,
  \end{align}
\end{subequations}
then matching coefficients of $\D S$ and $\D V$ tells us
\begin{equation}
T = \left(\frac{\partial e}{\partial S}\right)_{V},\quad\mbox{and}\quad
-p = \left(\frac{\partial e}{\partial V}\right)_{S}.
\end{equation}
Then we can relate $T$ and $p$ by taking the derivative of the $T$
equation with respect to $V$ while holding $S$ constant:
\begin{equation}
\left(\frac{\partial T}{\partial V}\right)_{S}
=\left(\frac{\partial}{\partial V}\left(\frac{\partial e}{\partial S}\right)_{V}
\right)_{S}
\end{equation}
and similarly differentiating $p$ with respect to $S$ while holding
$V$ constant:
\begin{equation}
-\left(\frac{\partial p}{\partial S}\right)_{V}
=\left(\frac{\partial}{\partial S}\left(\frac{\partial e}{\partial V}\right)_{S}
\right)_{V}.
\end{equation}
We assert both are equal to $(\partial^{2}e/\partial S\partial V)_{S,V}$,
giving us
\begin{equation}
\boxed{\left(\frac{\partial T}{\partial V}\right)_{S}=-\left(\frac{\partial p}{\partial S}\right)_{V}}
\end{equation}
which is called the \define{Maxwell Relation for (internal) Energy}.
\end{node}

\begin{remark}
There are other Maxwell relations, by looking at equations of state,
then manipulating in a similar manner (examine their difference
equation, set coefficients of differential equal to each other, then
differentiate and set the derivatives equal to each other).
\end{remark}

\begin{example}[Enthalpy]
Recall, we have the enthalpy
\begin{equation}
h = e + \frac{p}{\rho}=e+pV,
\end{equation}
then
\begin{subequations}
  \begin{align}
\D h &= \D e + p\,\D V + V\,\D p\\
&= (T\,\D S - p\,\D V) + p\,\D V + V\,\D p.
  \end{align}
\end{subequations}
This gives us the \define{Equation of State for Enthalpy}:
\begin{equation}
\boxed{\D h = T\,\D S + V\,\D p.}
\end{equation}
That's the first step accomplished.

Now, we can then define enthalpy as a function of entropy $S$ and
pressure $h=h(S,p)$ which gives us the difference equation
\begin{equation}
\D h = \left(\frac{\partial h}{\partial S}\right)_{p}\D S
 \left(\frac{\partial h}{\partial p}\right)_{S}\D p.
\end{equation}
Then comparing coefficients
\begin{equation}
T = \left(\frac{\partial h}{\partial S}\right)_{p},\quad\mbox{and}\quad
V = \left(\frac{\partial h}{\partial p}\right)_{S},
\end{equation}
and then after some dreary calculations we obtain the \define{Maxwell Relation for Enthalpy}:
\begin{equation}
\boxed{\left(\frac{\partial T}{\partial p}\right)_{S} = \left(\frac{\partial V}{\partial S}\right)_{p}.}
\end{equation}
\end{example}

\begin{definition}
We define the \define{Helmholtz free energy} $f := e - TS$.
\end{definition}

\begin{example}
The equation of state for the Helmholtz free energy is
\begin{equation}
\D f = - S\,\D T - p\,\D V.
\end{equation}
Then we can write the Helmholtz free energy as a function of
temperature and [specific] volume $f=f(T,V)$. Then we write the
difference equation for $f$, match the coefficients of the difference
equation with the equation of state, and so on, we obtain the
\define{Maxwell Relation for Helmholtz Free Energy}:
\begin{equation}\label{eq:maxwell-relation-for-helmholtz-free-energy}
\left(\frac{\partial S}{\partial V}\right)_{T}
=\left(\frac{\partial p}{\partial T}\right)_{V}.
\end{equation}
\end{example}

\begin{definition}
We define the \define{Gibbs Free Energy} $g := h - TS$.
\end{definition}

\begin{example}
Omitting the intermediate steps, the Maxwell relation for the Gibbs
Free Energy is:
\begin{equation}\label{eq:maxwell-relation-for-gibbs-free-energy}
\left(\frac{\partial S}{\partial p}\right)_{T}
=-\left(\frac{\partial V}{\partial T}\right)_{p}.
\end{equation}
\end{example}

\begin{summary}
We have thus obtained the 4 Maxwell relations (and there are only 4):
\begin{enumerate}
\item (Internal) Energy: $\displaystyle{\left(\frac{\partial T}{\partial V}\right)_{S}=-\left(\frac{\partial p}{\partial S}\right)_{V}.}$
\item Enthalpy: $\displaystyle{\left(\frac{\partial T}{\partial p}\right)_{S} = \left(\frac{\partial V}{\partial S}\right)_{p}.}$
\item Helmholtz free energy: $\displaystyle{\left(\frac{\partial S}{\partial V}\right)_{T}=\left(\frac{\partial p}{\partial T}\right)_{V}.}$
\item Gibbs free energy: $\displaystyle{\left(\frac{\partial S}{\partial p}\right)_{T}
=-\left(\frac{\partial V}{\partial T}\right)_{p}.}$
\end{enumerate}
\end{summary}

\begin{remark}
``Back in the day'', there was a mnemonic called the ``Tizsou Square''
(Wikipedia calls it the ``Thermodynamic Square'' and attributes it to
Max Born) which allowed you to derive the Maxwell relations for each
Thermodynamic quantity.
\end{remark}